\section{Design Goals}
\label{sec:vaios_design_goals}

The design of VAIOS is guided by the requirements of real-time embedded systems, where predictable execution, low overhead, and reliability are essential. The primary objective is to provide a runtime environment that supports structured application development without introducing unnecessary complexity or performance penalties.

A key goal is deterministic behavior. The system is designed to ensure that task scheduling and execution remain predictable under varying workloads. This is achieved through a priority-based preemptive scheduler, bounded interrupt handling, and controlled use of synchronization primitives. By avoiding hidden or implicit behaviors, VAIOS allows developers to reason about system timing more effectively.

Efficiency and minimal overhead are also central to the design. Core operations such as context switching, task scheduling, and inter-task communication are implemented to minimize runtime cost. Hardware-assisted mechanisms are leveraged wherever possible to reduce software overhead, ensuring that the system remains suitable for time-critical control applications.

Simplicity is another important consideration. The system avoids overly complex abstractions and focuses on a small, well-defined set of features that are sufficient for most embedded use cases. This reduces the learning curve and makes the system easier to debug and maintain, while still providing the necessary functionality for building complete applications.

Configurability is achieved through compile-time parameters, allowing the system to be tailored to different requirements without modifying core components. This includes control over scheduling behavior, memory limits, and feature enablement. Such flexibility enables VAIOS to adapt to both resource-constrained environments and more feature-rich deployments.

Finally, the system is designed with modularity in mind. Different subsystems, including scheduling, memory management, communication, and storage, are organized in a way that allows them to evolve independently. This supports incremental development and makes it easier to extend the system as new requirements emerge.